% ======================================================================
% Chapter 18: Operator Products and Effective Vertices
% (book pages 599-647; PDF pages 622-673)
% ======================================================================
\chapter{Operator Products and Effective Vertices}

Our analysis of QCD in Chapter~17 was founded on the principle of asymptotic
freedom, which told us that strong interaction processes with large momentum
transfer might reliably be treated in weak-coupling perturbation theory.  So
far, however, we have made little use in QCD of the more powerful tools of
the renormalization group.  In this chapter, we will work out some
implications of the Callan-Symanzik equation in QCD.  We will see that
asymptotically free theories have their own characteristic scaling behavior,
with corrections in the form of anomalous powers of logarithms of the
momentum scale.  Though these corrections are generally weaker than those in
the scalar field theories studied in Chapter~13, they nevertheless have
important qualitative effects on the strong interactions.

We begin by considering the scaling law for mass terms in QCD, taking over
directly the formalism that we used to describe the mass term of $\phi^4$
theory in Sections~12.4 and~12.5.  Other applications, however, require a
more powerful theoretical tool, the \emph{operator product expansion}.
Section~18.3 introduces a general description of products of operators in
quantum field theory and explains how such operator products are constrained
by the Callan-Symanzik equation.  The last two sections use this tool to
develop a new viewpoint toward deep inelastic scattering and other hard
processes in QCD.

\section{Renormalization of the Quark Mass Parameter}

Up to this point, we have always assumed that quark masses are small enough
that they can be ignored in high-energy processes.  This is not always an
adequate assumption even for the light quarks $u$, $d$, $s$; for the heavier
quarks $c$, $b$, $t$, the masses can have very important effects.  However,
since isolated quarks do not exist, it is not possible to define the mass of a
quark unambiguously.  In the discussion to follow, we will consider the quark
mass to be a parameter of QCD perturbation theory, defined by a
renormalization prescription at some renormalization scale $M$.

Because we define the quark mass as we would a coupling constant, by a
renormalization convention, we should expect that this parameter will run
according to a renormalization group evolution, so that different values of
the mass parameter apply to different processes.  We say that our original
prescription leads to an \emph{effective quark mass}, which depends on the
momentum scale at which it is evaluated.  In this section, we will work out
the leading dependence of this effective mass on the momentum scale.

The basic formalism for effective mass terms was set out in Section~12.5.
To add a mass term to the QCD Lagrangian, we must first define the mass
operator $(\bar{q}q)$ by a renormalization prescription at a scale $M$.  Then
we can define the quark mass by adding to the Lagrangian the term
\begin{equation}
  \Delta\mathcal{L}_m = -m\,(\bar{q}q)_M .
  \label{eq:18.1}
\end{equation}
In this discussion, we will assume that the quark mass $m$ is small enough
that we need only keep terms of leading order in $m$.  We will also assume,
for simplicity, that we have such a mass term for only one quark flavor.

In the zero-mass limit, Green's functions of the operator $(\bar{q}q)$ with
quark fields,
\begin{equation}
  G^{(n;k)}(x_1, \ldots, x_n; y_1, \ldots, y_n; z_1, \ldots, z_k)
  = \bigl\langle q(x_1)\cdots q(x_n)\bar{q}(y_1)\cdots\bar{q}(y_n)\,
    \bar{q}q(z_1)\cdots\bar{q}q(z_k) \bigr\rangle ,
  \label{eq:18.2}
\end{equation}
obey the Callan-Symanzik equation
\begin{equation}
  \Bigl[ M\frac{\partial}{\partial M} + \beta\frac{\partial}{\partial g}
    + 2n\gamma + k\gamma_{\bar{q}q} \Bigr]\ G^{(n;k)}\bigl( \{x_i\}; g, M
    \bigr) = 0 ,
  \label{eq:18.3}
\end{equation}
where $\gamma$ is the anomalous dimension of the quark field and
$\gamma_{\bar{q}q}$ is the anomalous dimension of the operator
$\bar{q}q$.  If we include the mass terms in the Lagrangian according to
\eqref{eq:18.1}, the Green's function of $n$ quark fields and $n$ antiquark
fields satisfies
\begin{equation}
  \Bigl[ M\frac{\partial}{\partial M} + \beta\frac{\partial}{\partial g}
    + 2n\gamma + m\frac{\partial}{\partial m} \Bigr]\ G^{(n)}\bigl(
      \{x_i\}, \{y_i\}; g, m, M \bigr) = 0 .
  \label{eq:18.4}
\end{equation}
The derivative with respect to $m$ counts the number of times the mass
operator is used.  In Section~12.5, we traded the variable $m$, with the
dimensions of mass, for a dimensionless variable.  However, in QCD, it is
just as convenient to consider the dimensionful parameter $m$ as a coupling
constant.  The solution of the Callan-Symanzik equation will then contain a
running mass parameter $\bar{m}(Q)$, which depends on a typical momentum $Q$
of the Green's function.  This parameter is defined as the solution to a
renormalization group equation analogous to Eq.~\eqref{eq:12.126}.  For this
case, the equation is
\begin{equation}
  \frac{d}{d\log(Q/M)}\ \bar{m} = \gamma_{\bar{q}q}(\bar{g})\cdot\bar{m} ,
  \label{eq:18.5}
\end{equation}
with the initial condition
\begin{equation}
  \bar{m}(M) = m .
  \label{eq:18.6}
\end{equation}
The quantity $\bar{m}(Q)$ is the \emph{effective mass}, which should be used
to compute the mass effects on quark production or scattering processes with
the momentum transfer $Q$.

To compute $\bar{m}(Q)$ explicitly, we need to work out the anomalous
dimension of the mass operator $\gamma_{\bar{q}q}$.  This can be done as
explained in Section~12.4.  We define the normalization of the operator
explicitly by the prescription that the vertex function of $(\bar{q}q)$
between renormalized quark fields should satisfy
\begin{equation}
  \frac{}{p^2-q^2=(p+q)^2=-M^2} = 1 .
  \label{eq:18.7}
\end{equation}
To preserve \eqref{eq:18.7}, we will need a counterterm vertex $\delta_{\bar{q}q}$
with the structure of the operator insertion.  Then, as in Eq.~\eqref{eq:12.112},
the anomalous dimension is given to one-loop order by
\begin{equation}
  \gamma_{\bar{q}q} = \frac{\partial}{\partial\log M}\Bigl(
    -\delta_{\bar{q}q} + \delta_2 \Bigr) ,
  \label{eq:18.8}
\end{equation}
where $\delta_2$ is the counterterm for the quark field strength
renormalization, defined in Fig.~\ref{fig:16.8}.  Correlation functions of the
gauge-invariant operator $(\bar{q}q)$ are gauge invariant, and so the various
terms in the Callan-Symanzik equation for this function must sum to a
gauge-invariant result.  Since the leading coefficient of $\beta(g)$ is
independent of the gauge and of other conventions, it follows from
\eqref{eq:18.3} that the leading coefficient of $\gamma_{\bar{q}q}$ is also
convention independent.  The counterterms $\delta_2$ and $\delta_{\bar{q}q}$
both depend on the gauge.  This argument shows that the gauge dependence
must cancel in \eqref{eq:18.8}.  In the calculation to follow, and in the
other anomalous dimension calculations in this chapter, we will work
consistently in Feynman-'t Hooft gauge.

We have already computed the divergent part of the counterterm $\delta_2$ in
Feynman-'t Hooft gauge in Section~16.4.  Evaluating the group-theory factor
in the result \eqref{eq:16.77} for QCD, we find
\begin{equation}
  \delta_2 = \frac{4}{3}\ \frac{g^2}{(4\pi)^2}\ \frac{\Gamma(2-d/2)}
    {(M^2)^{2-d/2}} .
  \label{eq:18.9}
\end{equation}
To compute $\delta_{\bar{q}q}$, we must work out the one-loop correction to
the vertex \eqref{eq:18.7}.  This is given by the diagram
\begin{equation}
  \text{(one-loop correction to } \bar{q}q \text{ vertex)}
  = -ig^2\ \int \frac{d^dk}{(2\pi)^d}\ \frac{1}{(k+q)^2}\ \frac{1}{k^2}\
    \frac{1}{(k-p)^2}\ \cdots ,
  \label{eq:18.10}
\end{equation}
In the expression for this diagram, the factor 1 represents the $\bar{q}q$
operator insertion.  In the corresponding diagram for the renormalization of
the quark number current $j^\mu = \bar{q}\gamma^\mu q$, this factor would be
replaced by $\gamma^\mu$.  Since we need only the divergent part of the
vertex renormalization \eqref{eq:18.10}, we can approximate the integrand by
its value for large $k$.  Then this diagram becomes
\begin{equation}
  \text{(diagram)} \sim ig^2\ \frac{2}{(4\pi)^2}\ \Gamma(2-d/2)\ \cdots .
  \label{eq:18.11}
\end{equation}
To preserve the normalization condition \eqref{eq:18.7}, we must add the
counterterm
\begin{equation}
  \delta_{\bar{q}q} = \frac{8}{3}\ \frac{g^2}{(4\pi)^2}\ \frac{\Gamma(2-d/2)}
    {(M^2)^{2-d/2}} .
  \label{eq:18.12}
\end{equation}
Assembling \eqref{eq:18.8}, \eqref{eq:18.9}, and \eqref{eq:18.12}, we find
\begin{equation}
  \gamma_{\bar{q}q} = \frac{4}{3}\ \frac{g^2}{(4\pi)^2} .
  \label{eq:18.13}
\end{equation}

As we have noted in the previous paragraph, the anomalous dimension of the
quark number current can be found by a very similar calculation.  This will
give a good check on our formalism, since, as we have argued above
Eq.~\eqref{eq:12.110}, a conserved current is unambiguously normalized by its
integral, the conserved charge, and so must have zero anomalous dimension.
If we substitute $j^\mu$ for 1 in \eqref{eq:18.10} and use the same set of
approximations to reduce the integral, we find in the numerator the Dirac
matrix structure
\begin{equation}
  \gamma^\rho\gamma^\mu\gamma^\nu\gamma_\rho\gamma^\sigma\gamma_\mu
  \to \frac{1}{2}(-2)^2\,\gamma^\nu\,\cdots ,
  \label{eq:18.14}
\end{equation}
Then, instead of \eqref{eq:18.12}, we need the counterterm
\begin{equation}
  \delta_j = \frac{4}{3}\ \frac{g^2}{(4\pi)^2}\ \frac{\Gamma(2-d/2)}
    {(M^2)^{2-d/2}} .
  \label{eq:18.15}
\end{equation}
Combining this result with \eqref{eq:18.9}, we find
\begin{equation}
  \gamma_j = 0 ,
  \label{eq:18.16}
\end{equation}
in accord with our general arguments.

If we replace the gamma function in \eqref{eq:18.11} by an explicit factor
of $\log(\Lambda^2/Q^2)$, and then subtract the divergence using the
counterterm \eqref{eq:18.12}, we find that the vertex diagram behaves as
\begin{equation}
  \text{(vertex)} \sim \frac{4}{3}\ \frac{g^2}{8\pi^2}\
    \log\frac{M^2}{Q^2} .
  \label{eq:18.17}
\end{equation}
This diagram gives an enhancement at small external momenta.  Some of this
enhancement is associated with the (gauge-dependent) rescaling of the
external quark fields; relation \eqref{eq:18.8} tells us how to extract the
piece of this logarithm associated with the gauge-invariant enhancement of
the effective mass.  Thus, to order $\alpha_s$,
\begin{equation}
  \bar{m}(Q) = m\ \Bigl[ 1 + \frac{4}{3}\frac{\alpha_s}{\pi}
    \log\frac{M}{Q} + \cdots \Bigr] .
  \label{eq:18.18}
\end{equation}

\begin{figure}[t]
  \centering
  \includegraphics[width=0.55\textwidth]{fig18.1.png}
  \caption{Diagrams giving the leading logarithmic contributions to the
    momentum dependence of the quark effective mass.}
  \label{fig:18.1}
\end{figure}

To compute the momentum dependence of the effective mass more accurately,
we must take two more features of the calculation into account.  First, the
quantity $\alpha_s\log(M^2/Q^2)$ may become of order 1, and, in this case,
we must take into account all leading logarithmic terms of the form
$(\alpha_s\log(M^2/Q^2))^n$.  Contributions of this type come from all the
diagrams shown in Fig.~\ref{fig:18.1}.  Second, the coupling constant
$\alpha_s$ is itself a function of the momentum scale, giving a further
enhancement to contributions from small $Q$.  Both of these effects are
properly accounted by solving the renormalization group equation
\eqref{eq:18.5}.  To the leading order in $g^2$, this equation takes the
explicit form
\begin{equation}
  \frac{d}{d\log(Q/M)}\ \bar{m} = \frac{4}{3}\ \frac{\bar{g}^2}{16\pi^2}\,
    \bar{m} .
  \label{eq:18.19}
\end{equation}
Inserting the solution of the renormalization group equation for $g$ in the
form \eqref{eq:17.17}, we find
\begin{equation}
  \bar{m}(Q) = \bar{m}(\mu)\ \Bigl[ \frac{\log(\mu^2/\Lambda^2)}
    {\log(Q^2/\Lambda^2)} \Bigr]^{4/(3b_0)} .
  \label{eq:18.20}
\end{equation}
The effective quark mass decreases slowly as the momentum scale increases.
This is a specific, testable prediction of QCD: the running of the quark
masses.

\section{QCD Renormalization of the Weak Interactions}

In this section we will apply the renormalization group formalism to the weak
interactions.  The weak interactions of the quarks involve four-fermion
operators of the type we will discuss in detail in Chapter~20.  We will see
that the QCD corrections to these operators can be important, and that they
lead to operator mixing.

We recall that the weak interactions of the $u$, $d$, and $s$ quarks are
described by the effective Lagrangian
\begin{equation}
  \Delta\mathcal{L} = -\frac{G_F}{\sqrt{2}}\ \bigl( \bar{u}_L\gamma^\mu d_L
    \bigr)\ \bigl( \bar{e}_L\gamma_\mu\nu_{eL} \bigr)\ \cos\theta_c
    + \cdots ,
  \label{eq:18.28}
\end{equation}
where $\theta_c$ is the Cabibbo angle.  This term allows the $s$ quarks to
decay through the process $s\to uud$.  Similarly, the rotation of
\eqref{eq:18.28} produces the effective interaction
\begin{equation}
  \Delta\mathcal{L} = -\frac{4G_F}{\sqrt{2}}\ \sin\theta_c\,
    \bigl( \bar{s}_L\gamma^\mu u_L \bigr)\ \bigl( \bar{u}_L\gamma_\mu d_L
      \bigr) ,
  \label{eq:18.32}
\end{equation}
which leads to the decay $s\to u\bar{u}d$.  These weak interaction processes
are referred to as nonleptonic and semileptonic decay processes,
respectively.  Similar expressions apply to the other heavy quarks.

Given that Eq.~(18.31) and \eqref{eq:18.32} describe the weak
interaction coupling of the $s$ quark at a fundamental level, we now discuss
the modification of these couplings by QCD logarithms.  We have seen in the
previous section that QCD corrections have a profound effect in enhancing the
strength of the quark mass term of the underlying Lagrangian.  We will now
investigate whether the strength of the weak interactions can receive a
similar enhancement.

\begin{figure}[t]
  \centering
  \includegraphics[width=0.5\textwidth]{fig18.2.png}
  \caption{QCD correction to the strength of the semileptonic weak
    interaction vertex.}
  \label{fig:18.2}
\end{figure}

We first consider the semileptonic weak interaction operator \eqref{eq:18.32}.
The leptonic fermion bilinear is not affected by QCD, so the QCD enhancement
of this operator is just the same as that of its quark component
\begin{equation}
  \bar{u}_L\gamma^\mu s_L .
  \label{eq:18.33}
\end{equation}
However, this operator is a current and so has $\gamma = 0$.  In terms of
diagrams, the logarithmic enhancement resulting from the diagram shown in
Fig.~\ref{fig:18.2} is canceled by the quark field-strength renormalization,
as we saw already in our discussion of the current vertex in Section~18.1.
The left-handed projector $\tfrac{1}{2}(1-\gamma^5)$ commutes through the
diagram and has no effect on the final result.  The same remark applies to
the semileptonic weak interaction that links $u$ and $d$ quarks.  It implies,
for that case, that the normalization of the cross sections for deep inelastic
neutrino scattering given in \eqref{eq:17.35} is not affected by QCD
logarithms.

\begin{figure}[t]
  \centering
  \includegraphics[width=0.55\textwidth]{fig18.3.png}
  \caption{QCD corrections to the strength of the nonleptonic weak
    interaction vertex.}
  \label{fig:18.3}
\end{figure}

In the case of nonleptonic weak interactions, however, the effect of QCD is
not so simple.  Let us first compute the Feynman diagrams that give the
leading corrections to the renormalization of the weak interaction vertex
Eq.~(18.31) and then, at a later stage, build up the renormalization
group interpretation of these results.

At order $\alpha_s$, the nonleptonic weak interaction vertex receives
corrections from the diagrams shown in Fig.~\ref{fig:18.3}.  Notice that the
first diagram is precisely the current renormalization found in the
semileptonic case.  The second diagram gives the analogous renormalization
of the second quark current.  In the computation of $\gamma$, these two
contributions cancel the contributions from the field-strength
renormalization of the four quark fields.  The remaining four diagrams of
Fig.~\ref{fig:18.3} are new contributions which contribute potentially large
rescaling factors.

We now compute these diagrams, beginning with the third diagram in
Fig.~\ref{fig:18.3}.  As in the computation of Section~18.1, we are interested
in the logarithmically divergent contribution associated with values of the
loop momentum $k$ much larger than the external momenta.  The simplest way
to extract this contribution is to compute each diagram in the approximation
of zero external momentum.  In writing the expression for these diagrams, we
will omit the prefactor
\begin{equation}
  \frac{4G_F}{\sqrt{2}}\ \cos\theta_c\sin\theta_c .
  \label{eq:18.34}
\end{equation}
We will retain the quark fields to represent the external states, so that our
final expressions will have the form of rescaled operators.

Using this notation, the third diagram in Fig.~\ref{fig:18.3} has the value
\begin{equation}
  \text{(third diagram)} = ig^2\ \int \frac{d^dk}{(2\pi)^d}\ \frac{1}{k^4}\
    \bigl( \bar{d}_L\gamma^\rho\gamma^\lambda\gamma^\sigma t^a u_L \bigr)\
    \bigl( \bar{u}_L\gamma^\rho\gamma_\lambda\gamma_\sigma t^a s_L \bigr) .
  \label{eq:18.35}
\end{equation}
Using the symmetry of the $k$ integral, we extract the divergent piece:
\begin{equation}
  \begin{aligned}
    \text{(third diagram)} &= ig^2\ \int \frac{d^dk}{(2\pi)^d}\ \frac{1}{k^4}\
      \bigl( \bar{d}_L\gamma^\rho\gamma^\lambda\gamma^\sigma t^a u_L \bigr)\
      \bigl( \bar{u}_L\gamma^\rho\gamma_\lambda\gamma_\sigma t^a s_L \bigr)\\
    &= -g^2\ \frac{\Gamma(2-d/2)}{(4\pi)^2}\ \bigl( \bar{d}_L\gamma^\lambda
      \gamma^\mu\gamma_\lambda t^a u_L \bigr)\ \bigl( \bar{u}_L\gamma_\mu
        t^a s_L \bigr)\ \cdots .
  \end{aligned}
  \label{eq:18.36}
\end{equation}
To put the product of quark fields into a more familiar form, we apply the
Fierz transformation discussed at the end of Section~3.4.  If the color
matrices $t^a$ were not present, the product of fermion fields would be
exactly the one appearing in \eqref{eq:3.82}, and we would find
\begin{equation}
  \bigl( \bar{d}_L\gamma^\mu\gamma^\lambda\gamma_\mu u_L \bigr)\
    \bigl( \bar{u}_L\gamma_\lambda s_L \bigr) = 16\ \bar{d}_L\gamma^\mu u_L\
    \bar{u}_L\gamma_\mu s_L .
  \label{eq:18.37}
\end{equation}
The matrices $t^a$ redirect the color quantum numbers of the quark fields.
To clarify this, we need the analogue of identity \eqref{eq:3.77} for color.
To find this identity, consider the color invariant
\begin{equation}
  (t^a)_{ij}(t^a)_{k\ell} .
  \label{eq:18.38}
\end{equation}
The indices $i$, $k$ transform according to the 3 representation of color; the
indices $j$, $\ell$ transform according to the $\bar{3}$.  Thus,
\eqref{eq:18.38} must be a linear combination of the two possible ways to
contract these indices,
\begin{equation}
  A\,\delta_{i\ell}\delta_{kj} + B\,\delta_{ij}\delta_{k\ell} .
  \label{eq:18.39}
\end{equation}
The constants $A$ and $B$ can be determined by contracting \eqref{eq:18.38}
and \eqref{eq:18.39} with $\delta_{ij}$ and with $\delta_{jk}$ and adjusting
$A$ and $B$ so that the contractions of \eqref{eq:18.39} obey the identities
\begin{equation}
  \mathrm{tr}[t^a](t^a)_{k\ell} = 0;\qquad
  (t^at^a)_{i\ell} = \frac{4}{3}\delta_{i\ell} .
  \label{eq:18.40}
\end{equation}
This gives the identity
\begin{equation}
  (t^a)_{ij}(t^a)_{k\ell} = \frac{1}{2}\Bigl( \delta_{i\ell}\delta_{kj}
    - \frac{1}{3}\delta_{ij}\delta_{k\ell} \Bigr) .
  \label{eq:18.41}
\end{equation}
A similar relation holds for the generators of $SU(N)$ in the fundamental
representation, with $(1/3)$ replaced by $(1/N)$ in that case.

Inserting \eqref{eq:18.41} into \eqref{eq:18.36}, we find that the first
term of the identity generates a new four-fermion operator,
\begin{equation}
  \bigl( \bar{d}_{Li}\gamma^\mu u_{Lj} \bigr)\ \bigl( \bar{u}_{Lj}
    \gamma_\mu s_{Li} \bigr) ,
  \label{eq:18.42}
\end{equation}
where $i$, $j$ are color indices.  Applying the Fierz rearrangement in
\eqref{eq:18.37}, and then applying the additional rearrangement
\eqref{eq:3.79}, we can convert this operator to the form
\begin{equation}
  16\ \bigl( \bar{d}_{Li}\gamma^\mu u_{Lj} \bigr)\ \bigl( \bar{u}_{Lj}
    \gamma_\mu s_{Li} \bigr) = 16\ \bigl( \bar{u}_{Lj}\gamma^\mu u_{Lj}
    \bigr)\ \bigl( \bar{d}_{Li}\gamma_\mu s_{Li} \bigr) .
  \label{eq:18.43}
\end{equation}
The minus sign in \eqref{eq:3.79} is compensated by a minus sign from
interchanging the order of fermion fields.  The final result is a product of
color-singlet quark currents; however, the fields in these currents are
associated differently from the original operator.

The final result of our evaluation of this diagram is
\begin{equation}
  \text{(third diagram)} = -4g^2\ \frac{\Gamma(2-d/2)}{(4\pi)^2}\ \Bigl[
    \bar{d}_L\gamma^\mu u_L\ \bar{u}_L\gamma_\mu s_L
    - \frac{1}{3}\bar{d}_L\gamma^\mu u_L\ \bar{u}_L\gamma_\mu s_L
    \Bigr] .
  \label{eq:18.44}
\end{equation}
The fourth diagram of Fig.~\ref{fig:18.3} gives precisely the same
contribution.

The evaluation of the last two diagrams in Fig.~\ref{fig:18.3} is quite
similar.  The fifth diagram gives
\begin{equation}
  \text{(fifth diagram)} = ig^2\ \int \frac{d^dk}{(2\pi)^d}\ \frac{1}{k^4}\
    \bigl( \bar{d}_L\gamma^\rho\gamma^\lambda\gamma^\sigma t^a u_L \bigr)\
    \bigl( \bar{u}_L\gamma_\rho\gamma_\lambda\gamma_\sigma t^a s_L \bigr)\
    \cdots = +g^2\ \frac{\Gamma(2-d/2)}{(4\pi)^2}\ \bigl( \bar{d}_L\gamma^\mu
    t^a u_L \bigr)\ \bigl( \bar{u}_L\gamma_\mu t^a s_L \bigr) .
  \label{eq:18.45}
\end{equation}
The four-fermion operator can be simplified as follows, by the use of the
Fierz identity \eqref{eq:3.79}:
\begin{equation}
  \begin{aligned}
    \bigl( \bar{d}_L\gamma^\mu t^a u_L \bigr)\ \bigl( \bar{u}_L\gamma_\mu
      t^a s_L \bigr) &= +\bigl( \bar{s}_L\gamma^\mu\gamma^\lambda\gamma_\mu
      t^a s_L \bigr)\ \bigl( \bar{u}_L\gamma_\lambda t^a u_L \bigr)\\
    &= +(-2)^2\ \bigl( \bar{d}_L\gamma^\mu t^a u_L \bigr)\ \bigl(
      \bar{u}_L\gamma_\mu t^a s_L \bigr) .
  \end{aligned}
  \label{eq:18.46}
\end{equation}
Again, we must reduce the product of color matrices using identity
\eqref{eq:18.41}, and, again, the first term of this identity will require an
additional Fierz transformation.  The final result is
\begin{equation}
  \text{(fifth diagram)} = +g^2\ \frac{\Gamma(2-d/2)}{(4\pi)^2}\ \Bigl[
    \bar{d}_L\gamma^\mu u_L\ \bar{u}_L\gamma_\mu s_L
    - \frac{1}{3}\bar{d}_L\gamma^\mu u_L\ \bar{u}_L\gamma_\mu s_L
    \Bigr] .
  \label{eq:18.47}
\end{equation}
The last diagram in Fig.~\ref{fig:18.3} gives an identical contribution.
The sum of the contributions from these four diagrams is
\begin{equation}
  -3g^2\ \frac{\Gamma(2-d/2)}{(4\pi)^2}\ \Bigl[ \bar{d}_L\gamma^\mu u_L\
    \bar{u}_L\gamma_\mu s_L - \frac{1}{3}\bar{d}_L\gamma^\mu u_L\
      \bar{u}_L\gamma_\mu s_L \Bigr] .
  \label{eq:18.48}
\end{equation}

The extraction of the ultraviolet-divergent pieces of the diagrams of
Fig.~\ref{fig:18.3} is part of our formal prescription for computing the
Callan-Symanzik $\gamma$ function of the weak interaction vertex.  However, it
is useful to pause at this point and ask about the physical significance of
this divergence.  The diagrams of Fig.~\ref{fig:18.3} would not be divergent
if we computed them in the underlying theory with $W$ bosons.  In writing
the weak interaction as an effective local vertex, we approximated the $W$
boson propagator by a constant, assuming that the momentum $k$ that it
carried was much less than $m_W$:
\begin{equation}
  \frac{1}{k^2-m_W^2} \to -\frac{1}{m_W^2}\ \cdots .
  \label{eq:18.49}
\end{equation}
The approximation we used to compute the QCD corrections to the effective
vertex is valid only in the region of integration where $k^2 \ll m_W^2$.
Outside this region we must use the full $W$ propagator; this introduces an
extra factor of $k^2$ in the denominator and makes the integral converge.
Thus, in a direct calculation of the QCD correction, the ultraviolet-divergent
terms in the evaluation of Fig.~\ref{fig:18.3} would be replaced by
logarithms cut off at $m_W$.  The lower limit of the logarithm is set by the
external momenta.  In the decay of a $K$ meson---the lightest hadron
containing the $s$ quark---these are of order $m_K$.  Thus the correction
given in \eqref{eq:18.48} should be evaluated by replacing
\begin{equation}
  \frac{\Gamma(2-d/2)}{(4\pi)^2} \to \frac{1}{4\pi}\ \log\frac{m_W^2}{m_K^2}
  .
  \label{eq:18.50}
\end{equation}
With this interpretation, we can rewrite \eqref{eq:18.48} as the
order-$\alpha_s$ correction to the leading-order weak interaction vertex.  The
effect of this correction is the rescaling and modification of the weak
interaction operator:
\begin{equation}
  \begin{aligned}
    \bar{d}_L\gamma^\mu u_L\ \bar{u}_L\gamma_\mu s_L
    &\to \Bigl( 1 + \frac{\alpha_s}{4\pi}\log\frac{m_W^2}{m_K^2} \Bigr)\
      \bar{d}_L\gamma^\mu u_L\ \bar{u}_L\gamma_\mu s_L\\
    &\quad - 3\ \frac{\alpha_s}{4\pi}\log\frac{m_W^2}{m_K^2}\
      \bar{u}_L\gamma^\mu u_L\ \bar{d}_L\gamma_\mu s_L .
  \end{aligned}
  \label{eq:18.51}
\end{equation}
Notice that the QCD corrections not only rescale the normalization of the
original operator but also introduce a new operator with a different
structure.  This calculation makes concrete the idea introduced in
Section~12.4 that the diagrams that change the normalization of local
operators may also mix together different operators with the same dimension
and quantum numbers.

Since the value of the logarithm in \eqref{eq:18.50} is about 10, the size
of the leading QCD correction is of order 1 and so higher-order corrections
are important.  To sum the leading logarithmic corrections, we return to the
renormalization group analysis.  For clarity, define
\begin{equation}
  \mathcal{O}_1 = \bar{d}_L\gamma^\mu u_L\ \bar{u}_L\gamma_\mu s_L;
  \qquad \mathcal{O}_2 = \bar{u}_L\gamma^\mu u_L\ \bar{d}_L\gamma_\mu s_L .
  \label{eq:18.52}
\end{equation}
The anomalous dimension matrix for these operators can be read off from the
one-loop computation.  Solving the renormalization group equations, one finds
that the effective Hamiltonian for nonleptonic decays contains a
combination of these two operators with powers of the ratio of running
couplings, leading to a substantial enhancement of the $\Delta I = 1/2$
amplitude.  This effect, known as the enhancement of the $\Delta I = 1/2$
rule, is one of the classic successes of QCD.

\section{The Operator Product Expansion}

In the previous section we encountered an important example of a product of
local operators: the four-fermion operators of the weak interactions.  In this
section we will develop a general formalism for handling products of
operators in quantum field theory, the \emph{operator product expansion}
(OPE), which was proposed by Wilson.  This formalism will allow us to
analyze, in a systematic way, the behavior of matrix elements of products of
operators in the limit of short distances.

Consider a general Green's function of a product of two operators
$\mathcal{O}_1(x)$ and $\mathcal{O}_2(0)$, together with a set of fields
$\phi(y_i)$:
\begin{equation}
  G_{12}\bigl( x; y_1, \ldots, y_m \bigr)
  = \bigl\langle \mathcal{O}_1(x)\mathcal{O}_2(0)\ \phi(y_1)\cdots\phi(y_m)
    \bigr\rangle .
  \label{eq:18.63}
\end{equation}
For small $x$, this Green's function is singular, because the two operators
are at nearly the same point.  The singularities arise from the contractions
between the fields inside the two operators.  However, the singular
contractions can only involve a limited set of quantum numbers: the resulting
local operator must have the global symmetry quantum numbers of the product
of $\mathcal{O}_1\mathcal{O}_2$, but it is otherwise unrestricted.  It is useful
to write this operator as a linear combination of operators from a standard
basis.  The coefficients in this linear combination can depend on the
separation $x$.  Typically, products of operators in quantum field theory are
singular, so it is likely that some of the coefficients will have singularities
as $x\to 0$.  Combining these observations, Wilson proposed that the effects
of the operator product could be computed by replacing the product of
operators in \eqref{eq:18.63} with a linear combination of local operators,
\begin{equation}
  \mathcal{O}_1(x)\mathcal{O}_2(0) \to \sum_n C_{12n}(x)\,
    \mathcal{O}_n(0) ,
  \label{eq:18.64}
\end{equation}
where the coefficients $C_{12n}(x)$ are c-number functions.  This operator
product expansion (OPE) will depend only on the operators
$\mathcal{O}_1$, $\mathcal{O}_2$, and their separation and will be independent
of the identity and location of the other fields appearing in the Green's
function.

The expansion \eqref{eq:18.64} implies that the Green's function
\eqref{eq:18.63} can be expanded for small $x$ as follows:
\begin{equation}
  G_{12}\bigl( x; y_1, \ldots, y_m \bigr) = \sum_n C_{12n}(x)\,
    G_n\bigl( y_1, \ldots, y_m \bigr) ,
  \label{eq:18.65}
\end{equation}
where
\begin{equation}
  G_n\bigl( y_1, \ldots, y_m \bigr) = \bigl\langle \mathcal{O}_n(0)\,
    \phi(y_1)\cdots\phi(y_m) \bigr\rangle ,
  \label{eq:18.66}
\end{equation}
and all of the dependence on $x$ is now carried by the OPE coefficient
functions.  In the example of the previous section, the final amplitudes
depended in a rather involved way on the small separation of the two
operators, through the dependence of the coefficients in Eq.~(18.61) on
$m_W$.  From the viewpoint of the operator product expansion, this dependence
is carried by the coefficient functions and is determined for all matrix
elements when these are computed.

In Sections~18.1 and~18.2, we used the renormalization group to compute the
enhancement or suppression factors for operator matrix elements.  Thus it is
natural to expect that the form of the operator product coefficients is also
determined by the renormalization group.  We will now work out this relation.
To begin, we rewrite the expansion \eqref{eq:18.64} more precisely.  The
operators that appear in this relation must be defined at some renormalization
scale $M$.  Then the operator product expansion reads:
\begin{equation}
  \bigl[ \mathcal{O}_1(x) \bigr]_M \bigl[ \mathcal{O}_2(0) \bigr]_M
  = \sum_n C_{12n}(x; M)\ \bigl[ \mathcal{O}_n(0) \bigr]_M .
  \label{eq:18.67}
\end{equation}
Note that the coefficient functions can depend on $M$, since they must absorb
the $M$-dependent operator rescalings.  If we use the left-hand side of
\eqref{eq:18.67} to compute \eqref{eq:18.63}, this function obeys the
Callan-Symanzik equation
\begin{equation}
  \Bigl[ M\frac{\partial}{\partial M} + \beta\frac{\partial}{\partial g}
    + m\gamma + \gamma_1 + \gamma_2 \Bigr]\ G_{12}\bigl( x; y_1, \ldots,
      y_m; M \bigr) = 0 .
  \label{eq:18.68}
\end{equation}
Similarly, with the operator $\mathcal{O}_n$ normalized at $M$, the Green's
function \eqref{eq:18.66} obeys
\begin{equation}
  \Bigl[ M\frac{\partial}{\partial M} + \beta\frac{\partial}{\partial g}
    + m\gamma + \gamma_n \Bigr]\ G_n\bigl( y_1, \ldots, y_m; M \bigr) = 0 .
  \label{eq:18.69}
\end{equation}
By applying \eqref{eq:18.68} to the right-hand side of \eqref{eq:18.65}, we
see that these relations are consistent only if the OPE coefficient functions
obey the Callan-Symanzik equation
\begin{equation}
  \Bigl[ M\frac{\partial}{\partial M} + \beta\frac{\partial}{\partial g}
    + \gamma_1 + \gamma_2 - \gamma_n \Bigr]\ C_{12n}\bigl( x; M \bigr) = 0 .
  \label{eq:18.70}
\end{equation}
We now solve this equation by our standard methods.  First, let us apply
dimensional analysis.  If the operators $\mathcal{O}_1$, $\mathcal{O}_2$,
$\mathcal{O}_n$ have dimensions $d_1$, $d_2$, $d_n$, the coefficient function
$C_{12n}(x)$ must have the dimensions of $(\mathrm{mass})^{d_1+d_2-d_n}$.
Thus,
\begin{equation}
  C_{12n}(x) = \Bigl( \frac{1}{|x|} \Bigr)^{d_1+d_2-d_n}\
    \bar{C}(xM) ,
  \label{eq:18.71}
\end{equation}
where $\bar{C}(xM)$ is a dimensionless function.  This function is determined
from \eqref{eq:18.70} according to the method of Section~12.3.  Thus,
\begin{equation}
  C_{12n}(x) = \Bigl( \frac{1}{|x|} \Bigr)^{d_1+d_2-d_n}\ \bar{c}\bigl(
    \bar{g}(1/|x|) \bigr)\ \exp\Bigl[ \int_{1/|x|}^{M} d\log M'\ \bigl(
      \gamma_n - \gamma_1 - \gamma_2 \bigr) \Bigr] ,
  \label{eq:18.72}
\end{equation}
with $\bar{c}(g)$ a dimensionless function of the running coupling constant
at the separation scale $1/|x|$.

At a fixed point of the renormalization group, the $\gamma$ functions would
take definite values $\gamma_j^{*} = \gamma_j(g^{*})$.  Then, the solution
\eqref{eq:18.72} can be evaluated as
\begin{equation}
  C_{12n}(x) = \Bigl( \frac{1}{|x|} \Bigr)^{d_1+d_2-d_n}\ \bar{c}(g^{*})\
    \exp\Bigl[ \log(xM)\bigl( \gamma_n^{*} - \gamma_1^{*} - \gamma_2^{*}
      \bigr) \Bigr] .
  \label{eq:18.73}
\end{equation}
Thus, in this case,
\begin{equation}
  C_{12n}(x) = \bar{c}(g^{*})\ \Bigl( \frac{1}{|x|} \Bigr)^{d_1^{*}+d_2^{*}
    - d_n^{*}} ,
  \label{eq:18.74}
\end{equation}
where
\begin{equation}
  d_j^{*} = d_j + \gamma_j(g^{*})
  \label{eq:18.75}
\end{equation}
is the true scaling dimension of the operator $\mathcal{O}_j$ at the fixed
point.

For the case of an asymptotically free theory, the scaling relation is
complicated in the way that we worked out in Section~18.1.  In the leading
order of perturbation theory, the three $\gamma$ functions take the form
Eq.~(18.24).  Then the solution of \eqref{eq:18.72} takes the form
\begin{equation}
  C_{12n}(x) \sim \Bigl( \frac{1}{|x|} \Bigr)^{d_1+d_2-d_n}\ \Bigl(
    \frac{\log(1/|x|^2\Lambda^2)}{\log(M^2/\Lambda^2)} \Bigr)^{
      (\alpha_n-\alpha_1-\alpha_2)/2b_0} .
  \label{eq:18.76}
\end{equation}
In the example of Section~18.2, the original operators were currents with
dimension 3 and $\gamma = 0$, at separation $m_W^{-1}$, and the final local
operators had dimension 6.  Thus, \eqref{eq:18.76} does properly reproduce
the dependence of Eq.~(18.61) on $m_W$.  Notice that the renormalization
group dependence is less complicated for a product of currents, which have a
fixed normalization independent of scale.  This special case occurs often in
applications of the operator product expansion.

We have written Eq.~\eqref{eq:18.70} without taking account of operator
mixing.  However, as we have already seen, operator mixing is often an
essential part of the applications of the OPE.  It is straightforward to
include this effect by rewriting the analysis that leads to \eqref{eq:18.70}
using matrix-valued $\gamma$ functions.  For example, with operator mixing,
the Callan-Symanzik equation for $G_n$ will be modified to
\begin{equation}
  \Bigl[ M\frac{\partial}{\partial M} + \beta\frac{\partial}{\partial g}
    + m\gamma \Bigr]\ G_n + \sum_p \gamma_{np}\,G_p\bigl( y_1, \ldots, y_m; M
    \bigr) = 0 .
  \label{eq:18.77}
\end{equation}
With these changes, \eqref{eq:18.70} becomes
\begin{equation}
  \Bigl[ M\frac{\partial}{\partial M} + \beta\frac{\partial}{\partial g}
    \Bigr]\ C_{12n}(x; M) + \gamma_{1k}\,C_{12k}(x; M)
    + \gamma_{2k}\,C_{12n}(x; M) - C_{12k}(x; M)\,\gamma_{kn} = 0 .
  \label{eq:18.78}
\end{equation}
Notice that the first two $\gamma$ matrices act on the OPE coefficient from
the left, while the third acts from the right.  In the case of a product of
currents, the first two $\gamma$ matrices vanish and \eqref{eq:18.78}
simplifies to
\begin{equation}
  \Bigl[ M\frac{\partial}{\partial M} + \beta\frac{\partial}{\partial g}
    \Bigr]\ C_{12n}(x; M) - C_{12k}(x; M)\,\gamma_{kn} = 0 .
  \label{eq:18.79}
\end{equation}
This equation will play an important role in the analysis of Section~18.5.

\section{Operator Analysis of $e^{+}e^{-}$ Annihilation}

It is not difficult to imagine that there is a connection between matrix
elements in which currents are placed at short distances from one another and
matrix elements in which currents deliver a hard momentum transfer.  Thus we
might expect that the idea of the operator product expansion will give us a
new viewpoint from which to understand the theory of hard-scattering
processes in QCD.  In this section and the next, we will work out the relation
of the operator product expansion to the perturbative QCD analysis of
Chapter~17.

We begin by discussing the total cross section for $e^{+}e^{-}$ annihilation
to hadrons.  Below Eq.~\eqref{eq:17.9}, we argued that this total cross
section could be computed in QCD perturbation theory, using a value of
$\alpha_s$ corresponding to the scale of the total center of mass energy.
However, this argument was a purely intuitive one, with many logical jumps.
In this section, we will give a more rigorous argument to the same
conclusion.

\begin{figure}[t]
  \centering
  \includegraphics[width=0.5\textwidth]{fig18.4.png}
  \caption{Diagrams whose imaginary part yields the total cross section for
    $e^{+}e^{-}\to$ hadrons.}
  \label{fig:18.4}
\end{figure}

In order to invoke the operator product expansion, we must write the total
cross section for $e^{+}e^{-}$ annihilation to hadrons as the matrix element
of a product of currents.  To do this, we use the optical theorem to relate
the total $e^{+}e^{-}$ scattering cross section to the forward scattering
amplitude for $e^{+}e^{-}\to e^{+}e^{-}$.  Ignoring the mass of the electron,
we see from Eq.~\eqref{eq:7.49} that
\begin{equation}
  \sigma\bigl( e^{+}e^{-} \bigr) = \frac{1}{2s}\ \mathrm{Im}\,
    \mathcal{M}\bigl( e^{+}e^{-}\to e^{+}e^{-} \bigr) .
  \label{eq:18.80}
\end{equation}
To compute the cross section for $e^{+}e^{-}\to$ hadrons, we consider in the
computation of the imaginary part only the contributions from hadronic
intermediate states.  To leading order in $\alpha$, but to all orders in the
strong interactions, these contributions come from considering only diagrams
of the form of Fig.~\ref{fig:18.4}, and taking the imaginary part of the
hadronic contributions to the vacuum polarization.

The value of the diagrams shown in Fig.~\ref{fig:18.4} is
\begin{equation}
  i\mathcal{M} = (-ie)^2\ \bar{u}(k)\gamma^\mu v(k_+)\ \frac{1}{q^2}\
    \bigl( i\Pi^{\mu\nu}_H(q) \bigr)\ \frac{1}{q^2}\ \bar{v}(k_+)\gamma^\nu
    u(k) ,
  \label{eq:18.81}
\end{equation}
where $s = q^2$ and $\Pi^{\mu\nu}_H(q)$ is the hadronic part of the vacuum
polarization.  By the Ward identity, this can be written in the form
\begin{equation}
  \Pi^{\mu\nu}_H(q) = \bigl( g^{\mu\nu}q^2 - q^\mu q^\nu \bigr)\,
    \Pi_H(q^2) .
  \label{eq:18.82}
\end{equation}
The total hadronic cross section is then obtained by taking the imaginary
part:
\begin{equation}
  \sigma\bigl( e^{+}e^{-}\to\text{hadrons} \bigr) = \frac{4\pi\alpha^2}{s}\
    \mathrm{Im}\,\Pi_H(s) .
  \label{eq:18.83}
\end{equation}
Comparing with \eqref{eq:17.4}, we see that the function
$\mathrm{Im}\,\Pi_H(s)$ is the quantity that contains the QCD dynamics.

The vacuum polarization $\Pi_H(q^2)$ is a matrix element of a product of
electromagnetic currents,
\begin{equation}
  \Pi_H(q^2) \propto \int d^4x\ e^{iq\cdot x}\ \bigl\langle J^\mu(x)
    J^\nu(0) \bigr\rangle ,
\end{equation}
where the currents are the electromagnetic currents of the quarks.  For large
$q^2$, the integral is dominated by small $x$, and we can apply the operator
product expansion to the product of currents.  The OPE of two currents
contains the identity operator, the mass operator, and the operator
$(F^a_{\mu\nu})^2$, among others:
\begin{equation}
  J^\mu(x)J^\nu(0) \sim C_1(x)\,\mathbb{1} + C_{\bar{q}q}(x)\,
    m\,\bar{q}q(0) + C_{G^2}(x)\, (F^a_{\mu\nu})^2(0) + \cdots .
  \label{eq:18.91}
\end{equation}
Taking the vacuum expectation value, the leading term gives
\begin{equation}
  \Pi_H(q^2) \approx \Pi_H^{(0)}(q^2) + \mathcal{O}\bigl( 1/q^4 \bigr) ,
  \label{eq:18.94}
\end{equation}
where $\Pi_H^{(0)}$ is the leading perturbative contribution, computed in
QCD perturbation theory with the running coupling evaluated at the scale
$q^2$:
\begin{equation}
  \Pi_H^{(0)}(q^2) = \frac{1}{12\pi^2}\ \sum_f Q_f^2\ \Bigl( 1 +
    \frac{\alpha_s(q^2)}{\pi} + \cdots \Bigr) .
  \label{eq:18.95}
\end{equation}
This gives the total cross section for $e^{+}e^{-}$ annihilation, with the
same result that we derived in Chapter~17.

\begin{figure}[t]
  \centering
  \includegraphics[width=0.5\textwidth]{fig18.6.png}
  \caption{Analytic singularities of $\Pi_H(q^2)$ in the complex $q^2$
    plane.}
  \label{fig:18.6}
\end{figure}

\begin{figure}[t]
  \centering
  \includegraphics[width=0.55\textwidth]{fig18.7.png}
  \caption{Contour of integration involved in the derivation of the ITEP sum
    rules for $\sigma(e^{+}e^{-}\to$ hadrons).}
  \label{fig:18.7}
\end{figure}

The operator product expansion gives additional information beyond the
leading-order result, through the higher-dimension operators in
\eqref{eq:18.91}.  These corrections can be related to the observable cross
section by using dispersion relations.  Consider the integral
\begin{equation}
  \frac{1}{2\pi i}\oint dq^2\ \frac{1}{(q^2+Q_0^2)^{n+1}}\ \Pi_H(q^2) ,
  \label{eq:18.97}
\end{equation}
for $n > 1$, evaluated on a contour encircling $q^2 = -Q_0^2$.  If we
contract the contour onto the pole, we find
\begin{equation}
  \frac{1}{2\pi i}\oint dq^2\ \frac{1}{(q^2+Q_0^2)^{n+1}}\ \Pi_H(q^2)
  = \frac{1}{n!}\ \frac{d^n}{dq^{2n}}\ \Pi_H(q^2)\Big|_{q^2=-Q_0^2} ,
\end{equation}
which can be computed by evaluating $\Pi_H$ from the operator product
relation \eqref{eq:18.91},
\begin{equation}
  \Pi_H(q^2) = -e^2\ \Bigl[ c_1(q^2) + c_{\bar{q}q}(q^2)\,
    \langle 0|m\bar{q}q|0\rangle + c_{G^2}(q^2)\,
      \langle 0|(F^a_{\mu\nu})^2|0\rangle + \cdots \Bigr] .
  \label{eq:18.98}
\end{equation}
On the other hand, we can evaluate the integral by distorting the contour to
the form of Fig.~\ref{fig:18.7}.  Since none of the coefficient functions grow
faster than $(q^2)^0$ times logarithms as $q^2\to\infty$, the contour at
infinity can be neglected for $n > 1$.  The piece of the contour that wraps
around the branch cut gives
\begin{equation}
  \begin{aligned}
    &= -\frac{1}{4\pi}\ \int_0^{\infty} \frac{dq^2}{(q^2+Q_0^2)^{n+1}}\
      \mathrm{Disc}\,\Pi_H(q^2)\\
    &= \frac{1}{2\pi}\ \int_0^{\infty} \frac{ds}{(s+Q_0^2)^{n+1}}\
      \sigma(s) ,
  \end{aligned}
  \label{eq:18.99}
\end{equation}
where $\sigma(s)$ is the total cross section for $e^{+}e^{-}\to$ hadrons.
By equating \eqref{eq:18.97} and \eqref{eq:18.99}, we obtain a series of
integral relations between the OPE coefficients, evaluated in QCD perturbation
theory, and the observable cross section.  These relations, which were first
constructed by Novikov, Shifman, Voloshin, Vainshtein, and Zakharov, are
known as the \emph{ITEP sum rules}.%
\footnote{The theory of these sum rules is reviewed in V.~A.~Novikov,
L.~B.~Okun, M.~A.~Shifman, A.~I.~Vainshtein, M.~B.~Voloshin, and
V.~I.~Zakharov, \emph{Phys.\ Repts.} \textbf{41}, 1 (1978).}

Evaluating the sum rules with only the leading QCD expression for
$c_2(q^2)$, we find
\begin{equation}
  \int_0^{\infty} \frac{ds}{(s+Q_0^2)^{n+1}}\ \sigma(s)
  = \frac{2\pi\alpha^2}{3}\ \sum_f Q_f^2\ \frac{1}{Q_0^{2n}}\ \Bigl[ 1 +
    \mathcal{O}\bigl( \alpha_s(Q_0^2) \bigr) + \mathcal{O}\Bigl(
      \frac{\langle G^2\rangle}{Q_0^4} \Bigr) + \cdots \Bigr] .
  \label{eq:18.100}
\end{equation}
The leading-order relation is consistent with the lowest-order cross section
given in Eq.~\eqref{eq:18.95}.  The corrections come from higher orders of
QCD perturbation theory, with $\alpha_s$ taken at the scale $Q_0^2$, and
from the higher operator terms in the OPE.

If the correction terms in \eqref{eq:18.100} converged to zero uniformly in
$n$, we could invert the sum rules and derive from them our result
\eqref{eq:18.94}.  However, the true situation is more subtle.  Because the
derivatives in \eqref{eq:18.97} emphasize terms with stronger $q^2$
variation, the correction terms in the ITEP sum rules are more and more
important as $n$ increases.  Thus the most important deviations of the cross
section from the prediction of QCD perturbation theory are oscillations about
this prediction, which average out in the sum rules for low $n$.  The
comparison of theory and experiment is shown in Fig.~\ref{fig:18.8}.  At
large $s$, \eqref{eq:18.94} is quite accurate.  As $s$ becomes smaller,
however, the oscillations grow in size.  Eventually, they come to dominate
the total cross section as the resonances associated with quark-antiquark
bound states.

\begin{figure}[t]
  \centering
  \includegraphics[width=0.6\textwidth]{fig18.8.png}
  \caption{Experimental measurements of the total cross section for the
    reaction $e^{+}e^{-}\to$ hadrons at energies below 3 GeV, compared to the
    prediction of perturbative QCD for 3 quark flavors.  The data are taken
    from the compilation of M.~Swartz, \emph{Phys.\ Rev.\ D} \textbf{53},
    5268 (1996).  Complete references to the various results are given
    there.}
  \label{fig:18.8}
\end{figure}

\section{Operator Analysis of Deep Inelastic Scattering}

We now apply the operator product expansion to another example of a QCD
hard-scattering process, deep inelastic electron scattering.  In Chapter~17 we
found that the predictions of QCD for deep inelastic scattering are precise
but also intricate in structure.  At a first level, QCD implies that deep
inelastic scattering is described by the parton model, in which the incident
electron scatters from quarks and antiquarks that carry fractions of the total
momentum of the proton.  These fractions are determined by parton
distribution functions, which reflect the form of the proton wavefunction and
are determined by soft QCD dynamics.  However, we saw in Section~17.5 that
effects of QCD perturbation theory cause the parton distributions to change
their form as a function of the momentum transfer $Q^2$.  We will now show
that much of this picture can be reconstructed from our new viewpoint, using
the operator product expansion.

In the previous section, we derived the OPE relations for the $e^{+}e^{-}$
annihilation cross section in three steps.  First, we used the optical theorem
to relate this cross section to a matrix element of a product of currents.
Second, we applied the operator product expansion to the product of currents.
Unfortunately, this expansion could be used only in an unphysical kinematic
region.  However, in the third step, we used the method of dispersion
relations to connect this unphysical result to an integral over the cross
section we wished to predict.  In our discussion of deep inelastic scattering,
we will go through these same three steps.  To obtain our final result, we
will need to add a fourth step, involving QCD operator rescaling.

\subsection{Kinematics of Deep Inelastic Scattering}

We begin by writing a general expression for the deep inelastic scattering
cross section.  The matrix element for deep inelastic electron scattering to
a final state $f$ is computed as shown in Fig.~\ref{fig:18.9}(a):
\begin{equation}
  \mathcal{M} \propto \int d^4x\ e^{iq\cdot x}\ \langle f|J^\mu(x)|P\rangle ,
  \label{eq:18.101}
\end{equation}
where $J^\mu(x)$ is the quark electromagnetic current.  The core of this
expression is the hadronic matrix element of the current between the proton
and some high-energy hadronic state.  This matrix element must be squared
and summed over possible final states.  That sum can be computed, using the
optical theorem, by relating it to the forward matrix element of two currents
in the proton state, as shown in Fig.~\ref{fig:18.9}(b).  Define
\begin{equation}
  W^{\mu\nu}(P, q) = \int \frac{d^4x}{2\pi}\ e^{iq\cdot x}\
    \langle P|J^\mu(x)J^\nu(0)|P\rangle ,
  \label{eq:18.102}
\end{equation}
averaged over the spin of the proton.  This object is known as the
\emph{forward Compton amplitude}, since if it is evaluated at $q^2 = 0$ and
contracted with physical polarization vectors, it gives the forward amplitude
for photon-proton scattering:
\begin{equation}
  i\mathcal{M}(\gamma p \to \gamma p) = (ie)^2\ \epsilon^{*}_\mu(q)
    \epsilon_\nu(q)\ \bigl( -iW^{\mu\nu}(P, q) \bigr) .
  \label{eq:18.103}
\end{equation}
However, in the following discussion we will need to analyze \eqref{eq:18.102}
for general spacelike $q$ and for general polarization states.

\begin{figure}[t]
  \centering
  \includegraphics[width=0.6\textwidth]{fig18.9.png}
  \caption{Computation of the cross section for deep inelastic electron
    scattering: (a) general structure of the amplitudes; (b) application of
    the optical theorem.}
  \label{fig:18.9}
\end{figure}

The optical theorem for Compton scattering from a proton is
\begin{equation}
  \frac{1}{2}\sum_f \int d\Pi_f\ |\mathcal{M}(\gamma p \to f)|^2
  = 2\,\mathrm{Im}\,\mathcal{M}(\gamma p \to \gamma p) .
  \label{eq:18.104}
\end{equation}
In the generalization given in \eqref{eq:7.49}, this result extends to the
more general situation in which the initial and final photon polarizations
can differ arbitrarily.  Transcribing \eqref{eq:18.104} to $W^{\mu\nu}$, we
find
\begin{equation}
  2\,\mathrm{Im}\,W^{\mu\nu}(P, q) = \sum_f \int d\Pi_f\ \langle P|
    J^\mu(-q)|f\rangle\ \langle f|J^\nu(q)|P\rangle ,
  \label{eq:18.105}
\end{equation}
where $J^\mu(q)$ is the Fourier transform of the current.

We can now compute the deep inelastic cross section in terms of
$W^{\mu\nu}$, using \eqref{eq:18.105} to represent the square of the last
factor.  The cross section should be averaged over initial and summed over
final electron spins.  Thus,
\begin{equation}
  \sigma(ep \to eX) = \int \frac{d^3k'}{(2\pi)^3}\ \frac{1}{4kk'}\ \Bigl[
    \sum_{\text{spins}} \bar{u}(k)\gamma^\mu u(k')\ \bar{u}(k')\gamma^\nu
    u(k) \Bigr]\ \bigl( \frac{e^2}{q^2} \bigr)^2\ 2\,\mathrm{Im}\,W_{\mu\nu}
    (P, q) .
  \label{eq:18.106}
\end{equation}
The electron spinor product can be evaluated as
\begin{equation}
  \sum_{\text{spins}} \bar{u}(k)\gamma^\mu u(k')\ \bar{u}(k')\gamma^\nu
    u(k) = \frac{1}{2}\mathrm{tr}\bigl[ \slashed{k}\gamma^\mu
    \slashed{k}'\gamma^\nu \bigr]
  = 2\bigl( k^\mu k'^{\nu} + k^\nu k'^{\mu} - g^{\mu\nu}k\cdot k'
    \bigr) .
  \label{eq:18.107}
\end{equation}
It is useful to convert the integral over the final electron momentum $k'$
and scattering angle $\theta$ to an integral over the dimensionless variables
$x$ and $y$ that we introduced in Section~17.3.  These variables are given in
terms of the initial and final electron energies $k$ and $k'$ by
\begin{equation}
  x = \frac{Q^2}{2P\cdot q},\qquad
  y = \frac{P\cdot q}{P\cdot k} ,
  \label{eq:18.108}
\end{equation}
where $Q^2 = -q^2 = 2kk'(1-\cos\theta)$.

The general form of $W^{\mu\nu}$ is constrained by Lorentz invariance and
current conservation to be
\begin{equation}
  W^{\mu\nu}(P, q) = \Bigl( -g^{\mu\nu} + \frac{q^\mu q^\nu}{q^2} \Bigr)\,
    W_1(x, Q^2) + \Bigl( P^\mu - \frac{P\cdot q}{q^2}q^\mu \Bigr)\,
    \Bigl( P^\nu - \frac{P\cdot q}{q^2}q^\nu \Bigr)\,
    W_2(x, Q^2) ,
  \label{eq:18.113}
\end{equation}
where $W_1$ and $W_2$ are scalar structure functions.  The cross section for
deep inelastic scattering can be expressed in terms of these structure
functions.  In the parton model, $W_1$ and $W_2$ are given in terms of the
parton distribution functions.  The relation is
\begin{equation}
  W_2(x, Q^2) = \sum_f Q_f^2\,f_f(x, Q^2),\qquad
  2x\,W_1(x, Q^2) = W_2(x, Q^2) ,
  \label{eq:18.120}
\end{equation}
with the $y$ dependence characteristic of free fermions, as in
Eq.~\eqref{eq:17.125}.  Finally, substituting from \eqref{eq:18.120} for the
imaginary part of $W_2$, we recover this parton model expression precisely:
\begin{equation}
  \frac{d^2\sigma}{dx\,dy}\bigl( ep \to eX \bigr)
  = \frac{2\pi\alpha^2}{Q^4}\ \Bigl[ 1 + (1-y)^2 \Bigr]\
    \sum_f Q_f^2\,f_f(x) .
  \label{eq:18.123}
\end{equation}
This equation will give us a reference point for comparison with more general
expressions that we will derive as we continue our analysis.

\subsection{Expansion of the Operator Product}

Since the forward Compton amplitude is a matrix element of a product of
currents, an alternative strategy for calculating $W^{\mu\nu}$ is to expand
this product as a series of local operators.  Like the parton model
evaluation, this method makes use of asymptotic freedom.  However, in this
case, the assumption is applied more directly.  The computation of the
operator product coefficients will take place explicitly at a small distance
of order $1/Q$, and so we can calculate these coefficients in a perturbation
theory whose coupling constant is $\alpha_s(Q^2)$.

In the previous section, we computed the coefficients of operators that
contribute to the vacuum expectation value of the product of currents by
considering the various ways of contracting the quark fields in the product.
Here, we should note that the operator 1 does not contribute to the Compton
scattering amplitude.  The leading contributions come from operators that can
create and annihilate quarks in the proton wavefunction.

The most important terms in the operator product of two currents come from
products of two quark currents $\bar{q}_f\gamma^\mu q_f$ with quarks of the
same flavor.  Therefore we will begin by studying the OPE of the individual
quark currents.  To zeroth order in $\alpha_s$, the leading terms of the
operator product of quark currents are given by
\begin{equation}
  \bar{q}\gamma^\mu q(x)\ \bar{q}\gamma^\nu q(0)
  = \bar{q}(x)\gamma^\mu q(x)\ \bar{q}(0)\gamma^\nu q(0)
    + \text{(contractions)} ,
  \label{eq:18.124}
\end{equation}
where the contractions should be evaluated as Feynman propagators for the
quark fields.  The terms with explicit contractions are singular as
$x\to 0$; the remaining terms are nonsingular and thus less important in the
short-distance limit.  In the OPE of currents with quarks of different
flavor, there are no corresponding singular terms; we will argue below that
this conclusion is valid even beyond the leading order in $\alpha_s$.

To evaluate $W^{\mu\nu}$, we must take the Fourier transform of the terms in
\eqref{eq:18.124}, as indicated in \eqref{eq:18.102}.  When we do this, we
should remember that the propagators carry not only the Fourier transform
momentum $q$ but also whatever momentum is carried in through the quark
fields.  To take account of this, it is convenient to represent
\begin{equation}
  \int d^4x\ e^{iq\cdot x}\ \bar{q}(x)\gamma^\mu q(x)\ \cdots
  = \int d^4x\ e^{iq\cdot x}\ \bar{q}(0)\,e^{ix\cdot d}\,\gamma^\mu\,
    q(0)\ \cdots ,
  \label{eq:18.125}
\end{equation}
where the derivatives $d$ act to the right on the quark field.  Notice that
this contribution has the structure of the first diagram on the right in
Fig.~\ref{fig:18.10}.  Similarly, the second contraction indicated in
\eqref{eq:18.124} has the form of the second diagram in Fig.~\ref{fig:18.10}.

\begin{figure}[t]
  \centering
  \includegraphics[width=0.55\textwidth]{fig18.10.png}
  \caption{Expansion of the operator product of two currents in the
    short-distance limit.}
  \label{fig:18.10}
\end{figure}

In the short-distance limit, the momentum $q$ will be larger than any
external momentum entering the quark fields.  Thus we should expand
\begin{equation}
  \frac{1}{(d+q)^2} = \frac{1}{Q^2 - 2iq\cdot d + d^2}
  = \frac{1}{Q^2}\ \sum_{n=0}^{\infty}\ \Bigl( \frac{2iq\cdot d - d^2}
    {Q^2} \Bigr)^n .
  \label{eq:18.126}
\end{equation}
We will argue below that the terms with $d^2$ in the numerator are
unimportant and may be dropped.  However, we should retain all powers of the
ratio $(2iq\cdot d/Q^2)$.  This ratio has $Q^2$ in the denominator and so is
formally suppressed in the short-distance limit.  However, in the parton
model
\begin{equation}
  \frac{2P\cdot q}{Q^2} = \frac{1}{x} ,
  \label{eq:18.127}
\end{equation}
so, eventually, all of these terms must be equally important.  We will see
how this works in a moment.

The last step required to reduce the operator product \eqref{eq:18.124} to a
useful form is to reduce the product of Dirac matrices.  We know from
\eqref{eq:18.113} that, after we average over the proton spin,
$W^{\mu\nu}$ will be symmetric under the interchange of $\mu$ and $\nu$.
Thus, it does no harm to symmetrize the OPE.  We can then reduce the product
of three Dirac matrices to one by using the identity
\begin{equation}
  \frac{1}{2}\bigl( \gamma^\mu\gamma^\alpha\gamma^\nu + \gamma^\nu
    \gamma^\alpha\gamma^\mu \bigr) = g^{\mu\alpha}\gamma^\nu + g^{\nu\alpha}
    \gamma^\mu - g^{\mu\nu}\gamma^\alpha ,
  \label{eq:18.128}
\end{equation}
which is easily proved from the anticommutation relations.  By the use of
\eqref{eq:18.126} and \eqref{eq:18.128}, we can rewrite \eqref{eq:18.125} as
\begin{equation}
  \int d^4x\ e^{iq\cdot x}\ \bar{q}(x)\gamma^\mu q(x)\ \bar{q}(0)\gamma^\nu
    q(0) = \bar{q}(0)\ \Bigl[ \gamma^\mu(i d^\nu) + \gamma^\nu(i d^\mu)
      - g^{\mu\nu}\slashed{d} \Bigr]\ \sum_{n=0}^{\infty}\ \frac{(2iq\cdot d)^n}
        {Q^{2n+2}}\ q(0) ,
  \label{eq:18.129}
\end{equation}
We can remove the term $(\slashed{d})q$, which vanishes to leading order in
$\alpha_s$, since the quark field obeys the Dirac equation.  To compute $W_1$
and $W_2$, we can also drop the terms with explicit factors of $q^\mu$, since
these will eventually be organized into the general form \eqref{eq:18.113}.
Then, finally, \eqref{eq:18.125} takes the form
\begin{equation}
  \int d^4x\ e^{iq\cdot x}\ \bar{q}(x)\gamma^\mu q(x)\ \bar{q}(0)\gamma^\nu
    q(0) = -\frac{1}{Q^2}\ \sum_{n=0}^{\infty}\ \Bigl( \frac{2i}{Q^2}
      \Bigr)^n\ \bar{q}(0)\ \gamma^{(\mu}q^{\mu_1}\cdots q^{\mu_n)}
      \gamma^{\nu)}\ (id_{\mu_1})\cdots(id_{\mu_n})\,q(0) ,
  \label{eq:18.130}
\end{equation}
symmetrized under $\mu\leftrightarrow\nu$.

The second term in \eqref{eq:18.124} differs from the first by the
interchange of the points $x$ and 0 and the interchange of indices $\mu$ and
$\nu$.  Its Fourier transform is thus given by \eqref{eq:18.130} with the
replacement $q\to -q$.  The complete operator product therefore contains
only terms even in $q$.  All remaining contributions from the singular terms
of the operator product contain the operator
\begin{equation}
  \bar{q}\ (iD^{\mu_1})\cdots(iD^{\mu_k})\,q ,
  \label{eq:18.131}
\end{equation}
with an even number of indices, with these indices either identified with
$\mu$ or $\nu$ or contracted with powers of $q$.  To write the relevant terms
of the operator product expansion, we will modify this operator in two ways.
First, since the operator in \eqref{eq:18.131} has $n$ vector indices, it
contains components that transform under many different irreducible
representations of the Lorentz group.  Each component has a different
rescaling law under renormalization.  However, we will see below that only
the component of \eqref{eq:18.131} with the highest spin is relevant to our
analysis.  This component is obtained by totally symmetrizing the indices
$\mu_1, \ldots, \mu_n$ and then subtracting terms proportional to
$g_{\mu_i\mu_j}$ so that the operator is traceless on all pairs of indices.
We will retain only this component when we write out the operator product of
currents.  Second, the operator \eqref{eq:18.131} does not transform simply
under gauge transformations.  Since the original currents
$\bar{q}\gamma^\mu q$ were invariant to color gauge transformations, the
operator product of two currents must be a sum of gauge-invariant operators.
We can make \eqref{eq:18.131} gauge-invariant by replacing each factor of
$(id_{\mu})$ with a covariant derivative $(iD_\mu)$.  This modification adds
only terms proportional to the strong coupling constant $g$, so it has no
effect on our derivation of the operator product coefficients.

Incorporating these changes, let us define a spin-$n$ operator with quarks of
flavor $f$ as follows:
\begin{equation}
  \mathcal{O}^{f}_{\mu_1\cdots\mu_n} = \bar{q}_f\ \gamma^{(\mu_1}\
    (iD^{\mu_2})\cdots(iD^{\mu_n})\,q_f - \text{traces} ,
  \label{eq:18.132}
\end{equation}
with indices symmetrized and with appropriate subtractions.  We can use these
operators to write a final expression for the most singular part of the OPE
of two currents $J^\mu$.  The leading terms in this operator product come
from \eqref{eq:18.130} and the corresponding contraction with $q\to -q$.
Extracting the pieces of these expressions that contain the highest spin
operators \eqref{eq:18.132}, we find
\begin{equation}
  J^\mu(x)J^\nu(0) \sim \sum_f \sum_{n\ \text{even}} C_n(q^2)\,
    \mathcal{O}^{f\,\mu\nu}_{\mu_1\cdots\mu_n}\ q^{\mu_1}\cdots q^{\mu_n}
    + \cdots ,
  \label{eq:18.133}
\end{equation}
where the sums over $n$ run over even integers only.

Expression \eqref{eq:18.133} has been derived in the leading order in
$\alpha_s$.  Higher-order Feynman diagrams will contribute corrections to the
coefficient functions of order $\alpha_s(Q^2)$.  These corrections will be
important only if they are multiplied by large logarithms.  If we consider
the operators $\mathcal{O}^{(n)}$ appearing on the right-hand side to be
normalized at the renormalization scale $Q$, there is no large ratio of
momenta available to enhance the QCD corrections to the coefficient
functions.  Large logarithmic corrections may still arise at a later stage of
the calculation, when we compute the matrix elements of the operators
$\mathcal{O}^{(n)}$.

From the expansion \eqref{eq:18.133}, it is straightforward to compute an
expansion for $W^{\mu\nu}$ by taking its expectation value in the proton
state.  To carry out this computation, we need to know the proton matrix
elements of the operators $\mathcal{O}^{(n)}$.  Notice that these matrix
elements cannot depend on the direction of the momentum $q$, since that
dependence has been isolated in the coefficient functions.  This means that
only the proton momentum $P$ is available to carry the vector indices of the
matrix element.  We can therefore write the spin-averaged matrix element of
$\mathcal{O}^{(n)}$ as
\begin{equation}
  \langle P|\mathcal{O}^{(n)}_{\mu_1\cdots\mu_n}|P\rangle = A_n^f\cdot
    2\,P_{\mu_1}\cdots P_{\mu_n} - \text{traces} .
  \label{eq:18.134}
\end{equation}
The coefficients $A_n^f$ are dimensionless.  They are not quite pure numbers,
because they depend on the renormalization scale of the operators, but we
will treat them as constants in the next few paragraphs.

For the case $n = 1$, the operators $\mathcal{O}^{(1)}$ reduce simply to the
quark flavor currents $\bar{q}_f\gamma^\mu q_f$; in this case the operators
are normalized independently of any scale and the coefficients $A_1^f$ are
truly constants.  From our general discussion of form factors in
Section~6.2, we know that the proton matrix element of a conserved flavor
current at zero momentum transfer is given by
\begin{equation}
  \langle P|\bar{q}_f\gamma^\mu q_f|P\rangle = \bar{u}(P)\gamma^\mu u(P)\,
    F_{f1}(0) ,
  \label{eq:18.135}
\end{equation}
where $F_{f1}(0)$ is equal to the value of the corresponding conserved charge
in the proton state.  For the quark currents, this charge is just the number
of quarks (minus antiquarks) of flavor $f$ in the state $|P\rangle$, which
we will call $N_f$.  Averaging \eqref{eq:18.135} over the proton spin, we find
\begin{equation}
  \langle P|\bar{q}_f\gamma^\mu q_f|P\rangle = 2P^\mu\cdot N_f .
  \label{eq:18.136}
\end{equation}
Thus, for $n = 1$,
\begin{equation}
  A_1^f = N_f = \begin{cases}2 & u\\ 1 & d\\ 0 & \text{other}.\end{cases}
  \label{eq:18.137}
\end{equation}
Similarly, $\mathcal{O}^{(2)f}$ is the contribution of the quark flavor $f$
to the energy-momentum tensor of QCD:
\begin{equation}
  T^{\mu\nu} = \sum_f \mathcal{O}^{(2)f\,\mu\nu} + \cdots .
  \label{eq:18.138}
\end{equation}
Thus, $A_2^f$ is the fraction of the total energy-momentum of the proton that
is carried by the quark flavor $f$.

When we evaluate the series for $W^{\mu\nu}$ using \eqref{eq:18.133} and the
expression \eqref{eq:18.134} for the operator matrix elements, we find
\begin{equation}
  W^{\mu\nu} = \sum_f \Bigl[ \sum_{n\ \text{even}} C_n(q^2)\ A_n^f\ \cdots
    \Bigr] .
  \label{eq:18.139}
\end{equation}
Carrying out the sums, one recovers the parton model expression for the
structure functions, with the parton distribution functions given by the
appropriate moments of the operator matrix elements.

\subsection{The Moment Sum Rules}

The operator product expansion gives a systematic way to compute the
higher-twist corrections to the parton model, and to understand the
renormalization group evolution of the structure functions.  The key object
is the matrix element of the twist-2 operators.  These matrix elements obey
renormalization group equations, and the renormalization group evolution of
the moments of the structure functions follows from the anomalous dimensions
of the twist-2 operators.

The moment sum rules relate the moments of the parton distribution functions
to the operator matrix elements $A_n^f$.  Using dispersion relations, one
derives
\begin{equation}
  \int_0^1 dx\ x^{n-1}\,f_f^+(x, Q^2) = A_n^f ,
  \label{eq:18.154}
\end{equation}
for $n$ even.  These relations are known as the \emph{moment sum rules} for
the deep inelastic form factors.  They relate the $x$ moments of the parton
distribution functions, as defined by Eq.~(18.146), to the proton
matrix elements of twist-2 operators.

Because $W_2$ is a symmetric function of $\nu$, the moment sum rules apply
only for even $n$.  However, in deep inelastic neutrino scattering, there is
a third form factor in $W^{\mu\nu}$, associated with the interference term
between the vector and axial vector parts of the weak interaction current.
In Problem~18.2, we show that this form factor can be used to derive a set
of sum rules for odd $n$:
\begin{equation}
  \int_0^1 dx\ x^{n-1}\,f_f^-(x, Q^2) = A_n^f ,
  \label{eq:18.155}
\end{equation}
where $A_n^f$ is the coefficient of the proton matrix element
\eqref{eq:18.134} for odd $n$, and $f_f^-(x)$ is a form factor which, in the
parton model, evaluates to
\begin{equation}
  f_f^-(x) = f_f(x) - f_{\bar{f}}(x) .
  \label{eq:18.156}
\end{equation}
Combining this information with the argument given below \eqref{eq:18.136},
we can see that the definition of the parton distribution functions from the
deep inelastic form factors has the correct normalization.  Using
\eqref{eq:18.137}, we find
\begin{equation}
  \int_0^1 dx\ f_f^-(x) = N_f ,
  \label{eq:18.157}
\end{equation}
the (net) number of quarks of flavor $f$ in the proton.  Similarly,
\eqref{eq:18.154} and \eqref{eq:18.138} imply
\begin{equation}
  \int_0^1 dx\ x\,f_f^+(x) = \langle x\rangle_f ,
  \label{eq:18.158}
\end{equation}
where $\langle x\rangle_f$ is the fraction of the total energy-momentum of
the proton carried by quarks and antiquarks of flavor $f$.

\subsection{Operator Rescaling}

If the coefficients $A_n^f$ were truly constants, relations
\eqref{eq:18.154} and \eqref{eq:18.155} would be consistent with parton
distribution functions that satisfy Bjorken scaling.  However, as we remarked
below \eqref{eq:18.134}, these factors actually depend on $Q^2$, since this
is the normalization point of the operators in the operator product expansion
\eqref{eq:18.133}.  Since this dependence comes only through operator
rescaling, it involves only logarithms of $Q^2$, and so contributes only a
slow violation of Bjorken scaling.  We can work out the $Q^2$ dependence of
the parton distribution functions quantitatively by summing the leading
logarithmic corrections to the matrix elements of the twist-2 operators.

To account for these corrections, let us first assume (incorrectly, as we
will see below) that the twist-2 operators \eqref{eq:18.132} are renormalized
without operator mixing.  Then the leading logarithmic corrections to the
matrix element of the operator $\mathcal{O}^{(n)}$ would be summed by
rescaling the operator normalized at $Q$ to operators normalized at a
standard reference point $\mu$, of order 1 GeV.  The relation between these
conventions would be
\begin{equation}
  \bigl[ \mathcal{O}^{(n)} \bigr]_Q = \Bigl( \frac{\alpha_s(Q^2)}
    {\alpha_s(\mu^2)} \Bigr)^{a_n^{qq}/2b_0}\ \bigl[ \mathcal{O}^{(n)}
      \bigr]_\mu ,
  \label{eq:18.159}
\end{equation}
where $a_n^{qq}$ is the first coefficient of the $\gamma$ function of
$\mathcal{O}^{(n)}$.  Then the factors $A_n^f$ would depend on $Q^2$
according to
\begin{equation}
  A_n^f(Q^2) = \Bigl( \frac{\alpha_s(Q^2)}{\alpha_s(\mu^2)} \Bigr)^{
    a_n^{qq}/2b_0}\ A_n^f(\mu^2) .
  \label{eq:18.160}
\end{equation}
This equation agrees with the scale dependence of operator product
coefficients written in \eqref{eq:18.76}, for the special case of an operator
product of currents, $a_1 = a_2 = 0$.  To find the explicit form of the
rescaling factor, we must compute $a_n^{qq}$.

\begin{figure}[t]
  \centering
  \includegraphics[width=0.5\textwidth]{fig18.13.png}
  \caption{Diagrams contributing the anomalous dimension of the quark
    twist-2 operators.}
  \label{fig:18.13}
\end{figure}

To compute the $\gamma$ functions of the quark twist-2 operators, we must
compute their counterterms for operator rescaling.  These are determined by
the diagrams shown in Fig.~\ref{fig:18.13}.  It suffices to compute these
diagrams with external momentum $p$ entering through the quark line and zero
external momentum injected into the operator.  Under these conditions, the
matrix element of the operator $\mathcal{O}^{(n)}$, in leading order, equals
\begin{equation}
  \langle p|\mathcal{O}^{(n)}_{\mu_1\cdots\mu_n}|p\rangle
  = \bar{u}(p)\,\gamma_{\mu_1}\,p_{\mu_2}\cdots p_{\mu_n}\,u(p) .
  \label{eq:18.161}
\end{equation}
Here and at all later points in the discussion, we will treat the matrix
elements of $\mathcal{O}^{(n)}$ as though they are symmetrized in the $n$
indices and have all possible traces subtracted.  We must now evaluate the
diagrams of Fig.~\ref{fig:18.13} and collect all terms that rescale this
structure.

The first diagram of Fig.~\ref{fig:18.13} is quite straightforward to
evaluate.  We combine denominators using identity \eqref{eq:6.40}:
\begin{equation}
  \frac{1}{k^2(k-p)^2} = \int_0^1 dx\ \frac{1}{(k-xp)^4} ,
  \label{eq:18.163}
\end{equation}
the quantities in the denominator on the right are
$\bar{k} = k - xp$ and $\Delta = -x(1-x)p^2$.  We must now shift the
integral, substitute $\bar{k} = k + xp$ in the numerator, and pick out a term
proportional to $(n-1)$ powers of $p$.  If this term contains the factor
$g^{\mu_i\mu_j}$, we may drop it, since it contributes to the coefficient of
an operator of higher twist and since, in any event, it will be removed when
we subtract traces.  Thus, we must choose carefully which two factors of
$\bar{k}$ we replace with $\bar{k}$ when we replace the others with $(xp)$.
The following choices, simplified using the rotational symmetry of the
$\bar{k}$ integral, do not give useful contributions:
\begin{equation}
  \int \frac{d^4\bar{k}}{(2\pi)^4}\ \frac{\bar{k}^\mu\bar{k}^\nu}
    {(\bar{k}^2-\Delta)^4} = \frac{i}{(4\pi)^2}\ \frac{g^{\mu\nu}}{4}\
      \int d\bar{k}\ \frac{\bar{k}^2}{(\bar{k}^2-\Delta)^4} .
  \label{eq:18.164}
\end{equation}
The one remaining placement of the factors of $\bar{k}$ is
\begin{equation}
  \bar{k}^{\mu_1}\bar{k}^{\mu_n}\gamma^{\mu_1}\gamma^{\mu_n}
  \to \frac{1}{4}\bar{k}^2\,\gamma^{\mu_1}\gamma^{\mu_2}\ \cdots .
  \label{eq:18.165}
\end{equation}
Thus Eq.~(18.162) has the value
\begin{equation}
  \text{(first diagram)} = -\frac{2}{3}\,n(n+1)\ \frac{g^2}{(4\pi)^2}\
    \frac{1}{n(n+1)}\,\Gamma(2-d/2)\ \bar{u}(p)\,\gamma^{\mu_1}\,p^{\mu_2}
    \cdots p^{\mu_n}\,u(p)\ \cdots .
  \label{eq:18.166}
\end{equation}
It is not so obvious that there are additional contributions to the rescaling
of the operators $\mathcal{O}^{(n)}$.  Note, however, that the covariant
derivatives in \eqref{eq:18.132} contain explicit factors of the gauge field,
\begin{equation}
  iD_\mu = i\partial_\mu - gA^a_\mu t^a ,
  \label{eq:18.167}
\end{equation}
and these may be contracted with gauge field vertices on the external legs.
These contributions give rise to the second and third diagrams in
Figure~\ref{fig:18.13}.  The term in which two factors of $A^a_\mu$ from
\eqref{eq:18.167} are contracted with one another is proportional to
$G^{\mu\nu}$ and thus does not contribute to the rescaling of the
leading-twist operators.

The contributions we have just described have the form of sums over $j$,
where $\mu_j$ is the index of the derivative that includes the contraction.
Then the second diagram of Fig.~\ref{fig:18.13} is the sum over $j$ of the
following integral:
\begin{equation}
  \text{(second diagram)} = \sum_j\ \int \frac{d^4k}{(2\pi)^4}\ \frac{1}{k^2}
    \ \frac{1}{(k-p)^2}\ \bar{u}(p)\,\gamma^\mu\ \slashed{k}\ \gamma^{\mu_1}
    \cdots \cdots \,u(p) .
  \label{eq:18.168}
\end{equation}
Since $\mu_j$ and $\mu_1$ are symmetrized, we can use \eqref{eq:18.128} to
rewrite $\gamma^\mu\slashed{k}\gamma^{\mu_1} + \gamma^{\mu_1}\slashed{k}
\gamma^\mu$ in a convenient form.  Carrying out the integrals, one obtains
the anomalous dimension coefficients $a_n^{qq}$.  Combining these with the
contributions from the third diagram of Fig.~\ref{fig:18.13}, one finds the
full one-loop anomalous dimension matrix of the twist-2 operators.

The result of this computation is the anomalous dimension matrix
\begin{equation}
  \gamma^{(n)} = \frac{g^2}{8\pi^2}\ \begin{pmatrix}
    a_n^{qq} & a_n^{qg}\\ a_n^{gq} & a_n^{gg}
  \end{pmatrix} ,
  \label{eq:18.181}
\end{equation}
where the coefficients are given by
\begin{equation}
  \begin{aligned}
    a_n^{qq} &= \frac{4}{3}\Bigl[ 1 - \frac{2}{n(n+1)} + 4\sum_{j=2}^{n}
      \frac{1}{j} \Bigr],\\
    a_n^{qg} &= \frac{n^2+n+2}{n(n+1)},\\
    a_n^{gq} &= \frac{4}{3}\ \frac{n^2+n+2}{n(n^2-1)},\\
    a_n^{gg} &= 6\Bigl[ 1 - \frac{2}{n(n-1)} - \frac{2}{(n+1)(n+2)}
      + 4\sum_{j=2}^{n}\frac{1}{j} \Bigr] - \frac{2}{3}n_f .
  \end{aligned}
  \label{eq:18.181b}
\end{equation}
Notice that this matrix is not symmetric.  In the last line, $n_f$ is the
number of quark flavors, equal to 1 in this case; this term comes from
\eqref{eq:16.74}.

In the realistic case, QCD contains several quark flavors---$u$, $d$, $s$,
and also $c$ and $b$ when we work at momenta sufficiently large that we can
ignore the masses of these particles.  Then the anomalous dimension matrix
$\gamma^{(n)}$ has size $(n_f+1)\times(n_f+1)$.  The submatrix acting on
quark operators is diagonal, with all of the diagonal entries being given by
$a_n^{qq}$ in \eqref{eq:18.181}.  The quark-gluon and gluon-quark entries are
all given by $a_n^{qg}$ and $a_n^{gq}$, respectively, and are independent of
the flavor.  The gluon diagonal entry is given by $a_n^{gg}$ in
\eqref{eq:18.181} with the realistic value of $n_f$.  This means that the
gluon operator mixes with only one linear combination of quark operators:
\begin{equation}
  \sum_f \mathcal{O}^{(n)}_f ,
  \label{eq:18.182}
\end{equation}
the orthogonal linear combinations are simply rescaled, with the exponent
given by $a_n^{qq}$ or Eq.~(18.172).

Let us now apply this analysis of operator mixing to the evaluation of the
moment sum rules.  For odd $n$, there is no operator mixing, and so the $Q^2$
dependence of the right-hand side of \eqref{eq:18.155} is correctly given by
the simple rescaling \eqref{eq:18.160}.

For even $n$, we must take operator mixing into account.  The right-hand
side of the sum rule \eqref{eq:18.154} is the proton matrix element of a
twist-2 operator normalized at the scale $Q$.  Let us write an arbitrary
linear combination of these operators as
\begin{equation}
  c^i\ \bigl[ \mathcal{O}^{(n)}_i \bigr]_Q ,
  \label{eq:18.183}
\end{equation}
where the index $i$ runs over $g$ and the various flavors $f$.  To rescale
this operator to a fixed reference momentum $\mu$, we rewrite the
coefficients in a basis of left eigenvectors of $\gamma^{(n)}$ and rescale
each eigenvector according to \eqref{eq:18.159}.  In terms of the matrix
$a^n$ of rescaling coefficients, we can write the rescaling abstractly as
\begin{equation}
  c^i\ \bigl[ \mathcal{O}^{(n)}_i \bigr]_Q \to c^i\ \Bigl( \frac{\alpha_s(Q^2)}
    {\alpha_s(\mu^2)} \Bigr)^{a^n/2b_0}\ \bigl[ \mathcal{O}^{(n)}_i
      \bigr]_\mu .
  \label{eq:18.184}
\end{equation}
This rescaling, acting with $c^i$ to the left of the matrix $(a^n)$, is
precisely the prescription required by Eq.~\eqref{eq:18.79}.

Let us work this out explicitly for the case $n = 2$.  The right-hand side
of the moment sum rule \eqref{eq:18.154} is given by the matrix element of
$\mathcal{O}^{(2)}_f$.  We rewrite this as
\begin{equation}
  \sum_f \mathcal{O}^{(2)}_f = \Bigl( \sum_f \mathcal{O}^{(2)}_f \Bigr)\
    + \cdots ,
  \label{eq:18.185}
\end{equation}
where the first combination is the energy-momentum tensor, and the second
term mixes with the gluon operator $\mathcal{O}^{(2)}_g$.  The anomalous
dimension matrix acting on $(\mathcal{O}^{(2)}_f, \mathcal{O}^{(2)}_g)$ for
$n = 2$ has coefficients
\begin{equation}
  \begin{pmatrix}
    a_2^{qq} & n_f\,a_2^{qg}\\
    a_2^{gq} & a_2^{gg}
  \end{pmatrix}
  = \begin{pmatrix}
    \frac{16}{9} & \frac{2}{3}n_f\\
    \frac{16}{9} & -\frac{2}{3}n_f
  \end{pmatrix} .
  \label{eq:18.186}
\end{equation}
The left eigenvectors of this matrix, and their corresponding eigenvalues,
are
\begin{equation}
  (1, 1):\ a = 0;\qquad
  \Bigl( \frac{4}{3}, -\frac{n_f}{2} \Bigr):\ a = \frac{4}{3}\Bigl(
    \frac{16}{9} + \frac{2}{3}n_f \Bigr) .
  \label{eq:18.187}
\end{equation}
Notice that the first eigenvector gives a linear combination of operators
$\sum_f\mathcal{O}^{(2)}_f$ with zero anomalous dimension.  This operator is
in fact the total energy momentum tensor of QCD,
\begin{equation}
  T^{\mu\nu} = \sum_f \mathcal{O}^{(2)f\,\mu\nu} ,
  \label{eq:18.188}
\end{equation}
which must have $\gamma = 0$.  If we expand the second term in
\eqref{eq:18.185} in terms of the components \eqref{eq:18.187}, we can
compute the full form of the operator rescaling.  We find
\begin{equation}
  \bigl[ \mathcal{O}^{(2)}_f \bigr]_Q = \frac{1}{n_f}\ \Bigl[
    T^{\mu\nu} - \frac{4}{3}\frac{n_f}{16/3+n_f}\ \Bigl( \frac{\log(Q^2/
      \Lambda^2)}{\log(\mu^2/\Lambda^2)} \Bigr)^{-\lambda_2/2b_0}\ \cdots
    \Bigr] ,
  \label{eq:18.189}
\end{equation}
where $T^{\mu\nu}$ is the energy-momentum tensor \eqref{eq:18.188}.  The
right-hand side of the $n = 2$ moment sum rule is given by the coefficient
of the proton matrix element of this operator.  To evaluate this
coefficient, we need to define gluon analogues of the $A_n^f$, by writing,
analogously to \eqref{eq:18.134},
\begin{equation}
  \langle P|\mathcal{O}^{(n)g}_{\mu_1\cdots\mu_n}|P\rangle
  = A_n^g\cdot 2\,P_{\mu_1}\cdots P_{\mu_n} - \text{traces} .
  \label{eq:18.190}
\end{equation}
For the case $n = 2$, we note in particular that
\begin{equation}
  \langle P|T^{\mu\nu}|P\rangle = 2P^\mu P^\nu ,
  \label{eq:18.191}
\end{equation}
thus, \eqref{eq:18.188} implies
\begin{equation}
  \sum_f A_2^f + A_2^g = 1 .
  \label{eq:18.192}
\end{equation}
If we replace each operator in \eqref{eq:18.189} by the corresponding
coefficient $A_2^f(\mu)$, we will have an expression for the right-hand side
of the $n = 2$ moment sum rule which makes its $Q^2$ dependence explicit.

Although expression \eqref{eq:18.189} is rather complicated, it has a simple
form in the extreme limit $Q^2\to\infty$.  At asymptotic $Q^2$, the last two
terms of \eqref{eq:18.189} tend to zero, and the right-hand side of
\eqref{eq:18.189} becomes a fixed number times the energy-momentum tensor.
Then, using \eqref{eq:18.191}, we can evaluate the $n = 2$ moment sum rule
completely:
\begin{equation}
  \int_0^1 dx\ x\,f_f^+(x, \infty) = \frac{1}{n_f}\ .
  \label{eq:18.193}
\end{equation}
In this extreme limit, we find that each quark flavor carries the same fixed
fraction of the energy-momentum of the proton.  By \eqref{eq:18.192}, the
remainder is carried by the gluons.  To illustrate, in a theory with
$n_f = 4$, each quark flavor carries $3/28$ of the total momentum of the
proton, and the gluons carry the remaining $4/7$.
Figure~\ref{fig:18.16} shows how slowly these asymptotic results are
approached starting from realistic parton distributions.

\begin{figure}[t]
  \centering
  \includegraphics[width=0.6\textwidth]{fig18.16.png}
  \caption{Fractions of the total energy-momentum of the proton carried by
    various parton species, as a function of $Q$, according to the CTEQ fit
    to deep inelastic scattering data described in Fig.~\ref{fig:17.6}.  The
    $Q$ dependence of the curves is calculated from the QCD evolution
    equations.}
  \label{fig:18.16}
\end{figure}

\subsection{Relation to the Altarelli-Parisi Equations}

The operator mixing analysis just described gives predictions for the moments
of parton distributions which imply that these integrals are $Q^2$ dependent.
Of the various moment integrals that do not involve operator mixing, only
the $n = 1$ integrals which give the flavor quantum numbers of the proton are
constant as a function of $Q^2$.  The rest decrease as powers of $\log Q^2$.
Similarly, one linear combination of the matrix elements of $n = 2$ twist-2
operators remains constant with $Q^2$.  This relation is expressed by the sum
rule \eqref{eq:18.192}.  To write this relation more clearly, let us
introduce the parton distribution of gluons as a smooth function satisfying
the relations
\begin{equation}
  \int_0^1 dx\ x\,f_g(x, Q^2) = A_2^g(Q^2) .
  \label{eq:18.194}
\end{equation}
Then \eqref{eq:18.192} becomes just the total momentum sum rule for parton
distributions \eqref{eq:17.39}:
\begin{equation}
  \int_0^1 dx\ x\ \Bigl[ \sum_f f_f^+(x, Q^2) + f_g(x, Q^2) \Bigr] = 1 .
  \label{eq:18.195}
\end{equation}
It is not difficult to verify that, for $n > 2$, all of the eigenvalues of
the matrix $a^n$ of anomalous dimension coefficients are negative.  Thus, all
of the higher moment sum rules decrease, subject to the flavor charge and
momentum conservation laws.  In other words, the operator renormalization
analysis predicts that parton distributions shift down to smaller values of
$x$ as $\log Q^2$ increases.  It is pleasing that this is the same conclusion
that we reached in Section~17.5, where we derived the Altarelli-Parisi
equations to describe this evolution of the parton distributions.

\problems

\textbf{18.1} In this problem we will study the renormalization of the quark
mass parameter in QCD, following the analysis of Section~18.1.  Compute the
one-loop anomalous dimension $\gamma_{\bar{q}q}$ explicitly, verifying
\eqref{eq:18.13}, and check that the result is independent of the gauge
parameter.

\textbf{18.2} \emph{Sum rules for odd $n$.}  In the text we derived moment
sum rules for the deep inelastic structure functions that apply for even
values of $n$.  In this problem, show that the interference term between the
vector and axial vector parts of the weak interaction current in deep
inelastic neutrino scattering gives a set of sum rules for odd $n$, and
verify Eq.~\eqref{eq:18.155}.

\textbf{18.3} \emph{The anomalous dimension of the gluon twist-2 operator.}
Complete the computation of the one-loop anomalous dimension matrix of the
twist-2 operators, verifying the coefficients $a_n^{qg}$, $a_n^{gq}$, and
$a_n^{gg}$ in \eqref{eq:18.181}.

\textbf{18.4} In the analysis of the operator product expansion for the
weak interactions in Section~18.2, we found that QCD corrections mix the two
four-fermion operators.  In this problem, solve the renormalization group
equations for these operators and compute the effective enhancement factors
for the $\Delta I = 1/2$ and $\Delta I = 3/2$ amplitudes.
